\documentclass[11pt,a4paper,oneside]{book}
\usepackage{90-faelles/latex/studiearkiv}
\hypersetup{pdftitle={Matematik – Bind 3: Differentialligninger}}
\begin{document}
\frontmatter
\begin{titlepage}
\centering
\vspace*{30mm}
{\sffamily\color{arkivblaa}\Large MATEMATIK TIL UNIVERSITETSFYSIK\par}
\vspace{16mm}
{\Huge\bfseries Bind 3\\[6mm]Differentialligninger\par}
\vspace{13mm}
{\Large Formler, begreber og matematisk sammenhæng\par}
\vfill
\begin{minipage}{0.8\textwidth}
Ordinære og partielle differentialligninger, Greenfunktioner og spektrale randværdiproblemer.

Dette bind genbruger de samme LaTeX-kapitler som det samlede matematikkompendium. Planlagte kapitler er markeret som ufærdige. Der indgår ikke separate gennemregnede opgaveeksempler.
\end{minipage}
\vspace{14mm}
{\small Fysik-studiearkiv · Arbejdsudgave\par}
\end{titlepage}
\chapter*{Læsevejledning}
\addcontentsline{toc}{chapter}{Læsevejledning}
Dette bind er en emneopdelt udgave af det samlede matematikkompendium. Den generelle kilde og alle figurer vedligeholdes centralt under \texttt{00-matematik/chapters} og \texttt{00-matematik/figures}. For hele overblikket samt henvisninger mellem bindene anvendes den samlede udgave. Ufuldstændige kapitler er tydeligt markerede.
\begingroup\small\tableofcontents\endgroup
\mainmatter
\chapter{Ordinære differentialligninger}
\label{ch:ode}
En ordinær differentialligning forbinder en ukendt funktion af én variabel
med dens afledte. I fysik er den uafhængige variabel ofte tid, men kan
også være en rumkoordinat, en baneparameter eller en anden kontinuert
størrelse. Ligningen alene fastlægger sjældent en entydig løsning:
initialdata, definitionsområde og modellens regularitet er en del af
problemet.

Kapitlet behandler reelle skalarfunktioner og endeligdimensionale
systemer. Symbolet $t$ bruges for den uafhængige variabel, og en prik
eller prim betegner differentiation med hensyn til denne variabel.
Komplekse eksponentialfunktioner kan bruges som regneredskab, men reelle
fysiske løsninger skal sammensættes med de korrekte konjugerede bidrag.
Partielle differentialligninger behandles i det næste hovedemne.

\section{Eksistens, entydighed og begyndelsesbetingelser}
\label{sec:ode-eksistens}
\subsection{Orden, normalform og initialdata}
En ligning af orden $n$ kan på en regulær gren skrives
\[
 y^{(n)}=F(t,y,y',\ldots,y^{(n-1)}).
\]
En begyndelsesværdiformulering specificerer $y,y',\ldots,y^{(n-1)}$
i samme punkt $t_0$. Omskrivningen til normalform kan fejle, hvis man
skal dividere med en koefficient, der er nul. Sådanne punkter er
singulære for normalformen og skal undersøges særskilt.

Ved at indføre tilstandsvektoren
$\boldsymbol{x}=(y,y',\ldots,y^{(n-1)})^{\mathsf T}$ omskrives ligningen
til et førsteordenssystem
\begin{equation}
 \dot{\boldsymbol{x}}=\boldsymbol{F}(t,\boldsymbol{x}),\qquad
 \boldsymbol{x}(t_0)=\boldsymbol{x}_0.
 \label{eq:ode-ivp}
\end{equation}
Denne form samler eksistens- og stabilitetsteorien. Komponenter med
forskellige fysiske enheder bør skaleres, før en norm bruges som et mål
for en fysisk fejl. En matematisk tilstandsnorm kræver en fast valgt
skalering, ikke blot addition af størrelser med forskellige enheder.

\subsection{Lokale sætninger og maksimal fortsættelse}
\begin{samepage}
\begin{saetning}[Lokal eksistens og entydighed]
Lad $\boldsymbol{F}$ være kontinuert i en omegn af
$(t_0,\boldsymbol{x}_0)$ og lokalt Lipschitz i tilstanden, med en
Lipschitzkonstant som kan vælges ensartet i $t$ på et tilstrækkeligt
lille rektangel. Da har \eqref{eq:ode-ivp} en entydig lokal løsning.
Kontinuerte partielle afledte med hensyn til tilstanden er en
anvendelig tilstrækkelig betingelse for denne lokale Lipschitzegenskab.
\end{saetning}
\end{samepage}
Kontinuitet alene giver lokal eksistens via Peanos sætning i
endelig dimension, men ikke nødvendigvis entydighed. En formel løsning
eller en retningsfeltstegning beviser derfor ikke i sig selv, at
initialdata bestemmer én bane. Under entydighed kan to løsninger ikke
mødes i samme tilstand på samme tidspunkt og derefter skilles inden
for et fælles regulært område.

Den lokale løsning kan fortsættes til et maksimalt interval. Et endeligt
endepunkt kan skyldes, at tilstanden bliver ubegrænset, eller at banen
nærmer sig randen af modellens definitionsområde. Glathed af
$\boldsymbol{F}$ alene forhindrer ikke endeligt blow-up. En
lineær vækstgrænse
$\|\boldsymbol{F}(t,\boldsymbol{x})\|\le a(t)+b(t)\|\boldsymbol{x}\|$
med ikke-negative, lokalt integrable $a,b$ giver derimod kontrol over
væksten på endelige tidsintervaller, når modellen er defineret for alle
relevante tilstande.

\subsection{Integralform, iteration og følsomhed}
For kontinuert $\boldsymbol{F}$ er den differentiable løsning ækvivalent
med integralrelationen
\[
 \boldsymbol{x}(t)=\boldsymbol{x}_0+
 \int_{t_0}^{t}\boldsymbol{F}(s,\boldsymbol{x}(s))\dd s.
\]
Picard-iteration erstatter tilstanden under integralet med den foregående
approksimation. På et kort interval med $Lh<1$ er den tilsvarende
operator en kontraktion i supremumsnormen, når en passende invariant
funktionsmængde er valgt. Kravet om en kort tid er en del af argumentet;
det er ikke et krav om, at den faktiske løsning kun eksisterer kort tid.

Lad $\boldsymbol{z}$ være en differentiabel approksimation med residual
$\boldsymbol{r}=\dot{\boldsymbol{z}}-\boldsymbol{F}(t,\boldsymbol{z})$.
Hvis begge baner bliver i et område med Lipschitzkonstant $L$, giver
Grönwalls ulighed for $t\ge t_0$
\begin{align}
 \|\boldsymbol{z}(t)-\boldsymbol{x}(t)\|
 &\le\e^{L(t-t_0)}\|\boldsymbol{z}(t_0)-\boldsymbol{x}_0\|\notag\\
 &\quad+\int_{t_0}^{t}\e^{L(t-s)}\|\boldsymbol{r}(s)\|\dd s.
 \label{eq:ode-residual}
\end{align}
Estimater af residual og initialfejl har dermed en direkte betydning
for løsningsfejlen. Konstanten $L$ har dimension invers tid, når
$t$ er tid, så eksponenten er dimensionsløs. Grænsen kan være
konservativ; den er ikke en påstand om, at alle fejl faktisk vokser
med den viste eksponentialfunktion.

\section{Førsteordensligninger og integrerende faktor}
\label{sec:ode-foersteorden}
\subsection{Separation af variable og tabte løsninger}
For en separabel ligning $y'=a(t)b(y)$ kan man på et interval med
$b(y)\ne0$ definere en stamfunktion $H'(y)=1/b(y)$ og skrive
\[
 H(y(t))-H(y(t_0))=\int_{t_0}^{t}a(s)\dd s.
\]
Relationen kan være en fuldt brugbar implicit beskrivelse. En eksplicit
funktion $y(t)$ kræver desuden inversion af $H$ på en entydig gren.
Grænser og konvergens håndteres som i integralregningskapitlet.

Division med $b(y)$ fjerner mulige konstante løsninger med $b(y_*)=0$.
Disse skal undersøges før divisionen og føres tilbage til den endelige
løsningsbeskrivelse. At en implicit form indeholder en vilkårlig
konstant sikrer ikke, at alle ligevægte eller adskilte grene er med.
Under entydighed kan en ikke-stationær autonom løsning ikke krydse en
ligevægt i endelig tid i et fælles regulært område.

\begin{figure}[htbp]
\centering
\begin{tikzpicture}
\begin{axis}[width=.87\textwidth,height=61mm,axis lines=middle,
 xmin=0,xmax=3.2,ymin=-1.4,ymax=1.4,xlabel={$t$},ylabel={$y$},samples=100]
\foreach \tx in {.2,.7,1.2,1.7,2.2,2.7}{
 \foreach \yy in {-1,-.5,0,.5,1}{
  \pgfmathsetmacro{\dx}{.11/sqrt(1+(\yy)^2)}
  \pgfmathsetmacro{\dy}{-(\yy)*\dx}
  \pgfmathsetmacro{\xa}{\tx-\dx}\pgfmathsetmacro{\xb}{\tx+\dx}
  \pgfmathsetmacro{\ya}{\yy-\dy}\pgfmathsetmacro{\yb}{\yy+\dy}
  \edef\arrowcmd{\noexpand\draw[->,gray!65,thin]
   (axis cs:\xa,\ya)--(axis cs:\xb,\yb);}\arrowcmd
 }
}
\addplot[arkivblaa,thick,domain=0:3]{1.2*exp(-x)};
\addplot[arkivblaa,dashed,domain=0:3]{.6*exp(-x)};
\addplot[arkivrod,dashed,domain=0:3]{-.6*exp(-x)};
\addplot[arkivrod,thick,domain=0:3]{-1.2*exp(-x)};
\node[font=\small,anchor=north east] at (rel axis cs:.98,.98) {$y'=-y$};
\end{axis}
\end{tikzpicture}
\caption{Retningsfelt og analytiske baner for den dimensionsløse lineære
ligning $y'=-y$. Pilene viser hældning og stigende tid, mens kurverne
viser medlemmer af familien $y=C\e^{-t}$. Ligevægten $y=0$ bevares som
en selvstændig løsning og krydses ikke af de øvrige baner. Figuren er
en geometrisk illustration, ikke en trinvis opgaveløsning.}
\label{fig:ode-retningsfelt}
\end{figure}


\subsection{Lineære skalarligninger}
For kontinuerte $p,q$ på et interval er
$y'+p(t)y=q(t)$ lineær, også når koefficienterne varierer med tiden.
En normaliseret integrerende faktor og den initialiserede løsning er
\begin{align}
 \mu(t)&=\exp\!\left(\int_{t_0}^t p(s)\dd s\right),\notag\\
 y(t)&=\mu(t)^{-1}\left[y_0+
                    \int_{t_0}^t\mu(s)q(s)\dd s\right].
 \label{eq:ode-lineaer-scalar}
\end{align}
Produktreglen giver $(\mu y)'=\mu q$, hvilket forklarer konstruktionen.
Eksponentialfaktoren er forskellig fra nul, så der mistes ingen løsning
ved division med $\mu$. Funktionen $p$ har dimension invers tid, og
$q$ har samme dimension som $y'$.

Den homogene løsning beskriver fri udvikling fra initialdata; et
partikulært bidrag beskriver drivningen. Superposition gælder for en
lineær operator, men summen af to løsninger til samme inhomogene ligning
løser ikke automatisk den samme ligning: drivningerne summeres også.
Forskellen mellem to løsninger til samme drivning er homogen.

\subsection{Bernoulli, eksakte former og kendt partikulær løsning}
Følgende strukturer giver yderligere generelle transformationsregler.
Antagelser om ikke-nul størrelser og reelle potensgrene kontrolleres
før skiftet, ikke efter løsningen.
\begin{itemize}
\item \textbf{Bernoulli:} For $y'+py=qy^n$, hvor $n\ne0,1$, giver
$z=y^{1-n}$ på en positiv gren
$z'+(1-n)pz=(1-n)q$. For heltallige eller særligt valgte rationale
potenser kan andre grene også være mulige. En eventuel løsning $y=0$
skal undersøges i originalligningen; den kan gå tabt ved transformationen.
\item \textbf{Eksakt differentialform:} Formen
$M(t,y)\dd t+N(t,y)\dd y=0$ har et potentiale $\Psi$ med
$\Psi_t=M$, $\Psi_y=N$, hvis $M,N\in C^1$ på et simpelt sammenhængende
åbent område og $M_y=N_t$. Løsninger ligger på niveaukurver
$\Psi=C$. En graf $y(t)$ kræver lokalt $N\ne0$.
\item \textbf{Integrerende faktor til en form:} Hvis en ikke-nul faktor
$\rho$ gør $\partial_y(\rho M)=\partial_t(\rho N)$, bliver formen
eksakt. En faktor kun af $t$ kan søges, når
$(M_y-N_t)/N$ kun afhænger af $t$; da er
$\rho'/\rho=(M_y-N_t)/N$ på området med $N\ne0$.
\item \textbf{Riccati:} For $y'=a(t)y^2+b(t)y+c(t)$ og en kendt løsning
$y_p$ giver $y=y_p+1/u$ den lineære ligning
$u'+(2ay_p+b)u=-a$ på intervaller med $u\ne0$.
Den kendte løsning $y_p$ skal bevares som en særskilt mulighed.
\end{itemize}
En anden hyppig struktur er $y'=F(y/t)$ for $t\ne0$. Skiftet $v=y/t$
giver $t v'=F(v)-v$, som er separabel, hvor højresiden ikke er nul.
Betegnelsen homogen for denne skalastruktur må ikke forveksles med
betegnelsen homogen for en lineær ligning med drivning nul.

\section{Andenordensligninger, dæmpning og drivning}
\label{sec:ode-andenorden}
\subsection{Fundamentalsystem og Wronskideterminant}
For kontinuerte $p,q,r$ på $I$ er den normaliserede ligning
\[
 y''+p(t)y'+q(t)y=r(t).
\]
To homogene løsninger $y_1,y_2$ er et fundamentalsystem, hvis
Wronskideterminanten $W=y_1y_2'-y_1'y_2$ er forskellig fra nul i ét
punkt. Produktreglen og ligningen giver Abels identitet
\[
 W'(t)=-p(t)W(t),\qquad
 W(t)=W(t_0)\exp\!\left[-\int_{t_0}^{t}p(s)\dd s\right].
\]
Dermed forbliver $W$ forskellig fra nul på hele det sammenhængende
regulære interval. En nul-Wronskian er ikke i sig selv en universel
afhængighedstest for vilkårlige funktioner; her bruges afgørende, at
funktionerne løser samme regulære homogene ligning.

Den generelle løsning er $C_1y_1+C_2y_2+y_p$. Initialdata bestemmer
$C_1,C_2$ gennem en invertibel $2\times2$-matrix. Variation af konstanter
giver et partikulært bidrag
\begin{equation}
 y_p(t)=-y_1(t)\int_{t_0}^{t}\frac{y_2(s)r(s)}{W(s)}\dd s
        +y_2(t)\int_{t_0}^{t}\frac{y_1(s)r(s)}{W(s)}\dd s.
 \label{eq:ode-variation-scalar}
\end{equation}
De valgte nedre grænser giver $y_p(t_0)=y_p'(t_0)=0$. Andre
integrationskonstanter kan optages i den homogene del. Division med
Wronskianen er kun legitim på det regulære interval, hvor den er ikke-nul.

\subsection{Konstante koefficienter og multiplicitet}
For $a_2y''+a_1y'+a_0y=0$ med $a_2\ne0$ bestemmes de
karakteristiske rødder af $a_2\lambda^2+a_1\lambda+a_0=0$.
To forskellige rødder giver uafhængige eksponentialløsninger;
en dobbelt rod kræver både $\e^{\lambda t}$ og $t\e^{\lambda t}$.
For reelle koefficienter kommer ikke-reelle rødder som konjugerede par
og kan kombineres til reelle sinus-/cosinusløsninger.

Mere generelt giver en rod af multiplicitet $m$ bidragene
$t^j\e^{\lambda t}$ for $0\le j<m$ i en lineær ligning med konstante
koefficienter. Potenserne er ikke nye eksponentialrødder: de afspejler,
at operatorens nulrum er større end rummet af rene eksponentialfunktioner.
Koefficienten foran den højeste afledte skal stadig være ikke-nul.

\subsection{Den normaliserede dæmpede oscillator}
En nyttig fysisk normalform er
\begin{equation}
 y''+2\gamma y'+\omega_0^2y=f(t),\qquad
 \gamma\ge0,\quad\omega_0>0.
 \label{eq:ode-oscillator}
\end{equation}
Her har $\gamma,\omega_0$ dimension invers tid, og $f$ dimension
$[y]/[t]^2$. Hvis originalligningen har masse $m>0$, er $f$ kraften
divideret med $m$. Sæt $\tau=t-t_0$ og
$\omega_d=\sqrt{\omega_0^2-\gamma^2}$ i det underdæmpede tilfælde.
De homogene løsningsformer er
\begin{center}
\begin{tabular}{@{}ll@{}}
\toprule
Regime & $y_h(\tau)$\\
\midrule
$0\le\gamma<\omega_0$ & $\e^{-\gamma\tau}(C_1\cos\omega_d\tau+C_2\sin\omega_d\tau)$\\[2mm]
$\gamma=\omega_0$ & $\e^{-\gamma\tau}(C_1+C_2\tau)$\\[2mm]
$\gamma>\omega_0$ & $C_1\e^{r_+\tau}+C_2\e^{r_-\tau}$\\
\bottomrule
\end{tabular}
\end{center}
I sidste linje er $r_\pm=-\gamma\pm\sqrt{\gamma^2-\omega_0^2}$.
Ved positiv dæmpning går alle homogene bidrag mod nul fremad i tiden.
For $\gamma=0$ består oscillationen; dæmpning og forskydning af frekvens
er to forskellige virkninger. Figur~\ref{fig:ode-daempning} sammenligner
regimerne med samme skalerede initialdata.

\begin{figure}[htbp]
\centering
\begin{tikzpicture}
\begin{axis}[width=.87\textwidth,height=65mm,axis lines=middle,
 xmin=0,xmax=10,ymin=-1.15,ymax=1.2,xlabel={$\tau$},ylabel={$y$},
 legend style={font=\small,at={(.97,.97)},anchor=north east},samples=180]
\addplot[black,dashed,domain=0:10]{cos(deg(x))};\addlegendentry{$\gamma=0$}
\addplot[arkivblaa,thick,domain=0:10]{exp(-.4*x)*(cos(deg(sqrt(.84)*x))+.4/sqrt(.84)*sin(deg(sqrt(.84)*x)))};
\addlegendentry{$\gamma=0{,}4$}
\addplot[arkivrod,thick,domain=0:10]{exp(-x)*(1+x)};\addlegendentry{$\gamma=1$}
\addplot[black,dotted,very thick,domain=0:10]{exp(-1.5*x)*(cosh(sqrt(1.25)*x)+1.5/sqrt(1.25)*sinh(sqrt(1.25)*x))};
\addlegendentry{$\gamma=1{,}5$}
\end{axis}
\end{tikzpicture}
\caption{Analytiske homogene oscillatorresponser i dimensionsløse variable
med $\omega_0=1$ og samme initialdata $y(0)=1$, $y'(0)=0$.
Uden dæmpning består oscillationen. Under-, kritisk og overdæmpning
har forskellige tidsforløb; stærkere dæmpning betyder ikke automatisk
hurtigere sen tilnærmelse til ligevægt.}
\label{fig:ode-daempning}
\end{figure}


\subsection{Kausal respons og harmonisk drivning}
For nulinitialdata kan drivningens bidrag skrives
\[
 y_p(t)=\int_{t_0}^{t}g(t-s)f(s)\dd s,\qquad t\ge t_0.
\]
For $\tau>0$ er responskernen
\[
 g(\tau)=
 \begin{cases}
 \e^{-\gamma\tau}\sin(\omega_d\tau)/\omega_d,
     &\gamma<\omega_0,\\
 \tau\e^{-\gamma\tau},&\gamma=\omega_0,\\
 (\e^{r_+\tau}-\e^{r_-\tau})/(r_+-r_-),&\gamma>\omega_0.
 \end{cases}
\]
Kernen sættes til nul for $\tau<0$ og har $g(0)=0$, $g'(0^+)=1$.
Disse betingelser sikrer den korrekte normalisering: to afledte af
konvolutionen giver drivningen med koefficient 1. Kernens dimension er
tid. Den mere generelle distributionelle Greenfunktionsteori behandles
senere; her er integralet en almindelig kausal løsningsformel.

For en harmonisk drivning $f_0\cos(\Omega t)$ kan et stationært bidrag
repræsenteres som realdelen af $Y(\Omega)\e^{i\Omega t}$ med
\[
 Y(\Omega)=\frac{f_0}{\omega_0^2-\Omega^2+2i\gamma\Omega},\qquad
 |Y|=\frac{|f_0|}{\sqrt{(\omega_0^2-\Omega^2)^2+4\gamma^2\Omega^2}}.
\]
For $\gamma>0$ er dette endeligt, og den frie transient dør ud.
For forskydningsamplituden ligger en indre resonanstop ved
$\Omega=\sqrt{\omega_0^2-2\gamma^2}$, hvis
$0<\gamma<\omega_0/\sqrt2$; ellers ligger maksimum for
$\Omega\ge0$ ved nul. Resonansfrekvensen er derfor generelt hverken
$\omega_0$ eller den frie dæmpede frekvens $\omega_d$.

Ved $\gamma=0$ og $|\Omega|=\omega_0$ findes ikke et begrænset
stationært harmonisk bidrag af denne form. Den tidsafhængige resonante
respons må behandles med et sekulært led. At indsætte nul i nævneren
er ikke en beregning af en endelig fysisk amplitude.

\section{Lineære systemer og fundamentalløsninger}
\label{sec:ode-systemer}
For en kontinuert $n\times n$-matrix $A(t)$ og kontinuert drivning
$\boldsymbol{b}(t)$ er systemet
$\dot{\boldsymbol{x}}=A(t)\boldsymbol{x}+\boldsymbol{b}(t)$ lineært.
En normaliseret fundamentalmatrix opfylder
\[
 \Phi'(t)=A(t)\Phi(t),\qquad \Phi(t_0)=I.
\]
Kolonnerne er uafhængige homogene løsninger. Liouvilles formel giver
\[
 \det\Phi(t)=\exp\!\left(\int_{t_0}^{t}\operatorname{tr}A(s)\dd s\right),
\]
så $\Phi$ er invertibel på koefficienternes regulære interval. Systemets
initialiserede løsning er
\begin{equation}
 \boldsymbol{x}(t)=\Phi(t)\left[\boldsymbol{x}_0+
                  \int_{t_0}^{t}\Phi(s)^{-1}\boldsymbol{b}(s)\dd s\right].
 \label{eq:ode-variation-matrix}
\end{equation}
Formlen følger ved substitutionen $\boldsymbol{x}=\Phi\boldsymbol{c}$:
produktreglen eliminerer det homogene bidrag og efterlader
$\boldsymbol{c}'=\Phi^{-1}\boldsymbol{b}$. Matrixrækkefølgen må ikke
byttes om. Propagatoren $U(t,s)=\Phi(t)\Phi(s)^{-1}$ opfylder
$U(t,s)U(s,r)=U(t,r)$ og $U(s,s)=I$.

\subsection{Matrixeksponentialet og variable koefficienter}
For konstant $A$ er $\Phi(t)=\e^{A(t-t_0)}$, hvor
\[
 \e^{A\tau}=\sum_{j=0}^{\infty}\frac{\tau^jA^j}{j!}.
\]
Hvis $A=V\Lambda V^{-1}$, er $\e^{A\tau}=V\e^{\Lambda\tau}V^{-1}$.
Hvis $A$ ikke er diagonaliserbar, giver Jordanblokke yderligere
polynomier i $\tau$ gange eksponentialfunktioner. En gentagen egenværdi
alene betyder ikke, at en matrix mangler egenvektorer; det skal
undersøges via egenrummenes dimensioner.

For tidsafhængige matricer er
$\exp(\int A(t)\dd t)$ generelt ikke fundamentalløsningen, fordi matricer
på forskellige tidspunkter ikke behøver kommutere. Parvis kommutation
$A(t)A(s)=A(s)A(t)$ er en tilstrækkelig betingelse for denne enklere
form med faste integrationsgrænser. Ellers bruges fundamentalmatrix,
integralligning eller en tidsordnet konstruktion.

\section{Stabilitet, faseportrætter og perturbation}
\label{sec:ode-stabilitet}
\subsection{Autonome systemer og linearisering}
Et autonomt system har formen
$\dot{\boldsymbol{x}}=\boldsymbol{F}(\boldsymbol{x})$.
Et ligevægtspunkt $\boldsymbol{x}_*$ opfylder
$\boldsymbol{F}(\boldsymbol{x}_*)=0$. For $C^1$-funktioner fås den
lokale udvikling
\[
 \dot{\boldsymbol{\eta}}=J\boldsymbol{\eta}+o(\|\boldsymbol{\eta}\|),
 \qquad
 J=D\boldsymbol{F}(\boldsymbol{x}_*),\quad
 \boldsymbol{\eta}=\boldsymbol{x}-\boldsymbol{x}_*.
\]
Hvis alle egenværdier af $J$ har negativ realdel, er ligevægten lokalt
asymptotisk stabil. Hvis mindst én har positiv realdel, er ligevægten
ustabil. Når ingen har positiv realdel, men mindst én har realdel nul,
er lineariseringen generelt uafgjort for det ikke-lineære system.
For en skalar ligning reduceres de første to tests til fortegnet af
$F'(y_*)$; tilfældet $F'=0$ kræver yderligere analyse.

Et faseportræt viser tilstande og deres tidsretning, ikke direkte en
graf over tid. For en andenordensligning er de naturlige koordinater
$y$ og $y'$, eventuelt med en valgt skalering. Figur~\ref{fig:ode-fase}
viser tre forskellige lineære geometriske typer. Tiltrækning, rotation
og frastødning er uafhængige egenskaber af feltet.

\begin{figure}[htbp]
\centering
\begin{tikzpicture}
\begin{axis}[name=nodeplot,width=.31\textwidth,height=47mm,
 xmin=-2,xmax=2,ymin=-2,ymax=2,axis lines=middle,axis equal image,
 xlabel={$x_1$},ylabel={$x_2$},title={Stabil knude},title style={font=\small},
 tick label style={font=\scriptsize},samples=90]
\addplot[arkivblaa,thick,domain=0:2.5] ({1.7*exp(-x)},{1.7*exp(-2*x)});
\addplot[arkivblaa,thick,domain=0:2.5] ({-1.7*exp(-x)},{1.7*exp(-2*x)});
\addplot[arkivblaa,thick,domain=0:2.5] ({1.7*exp(-x)},{-1.7*exp(-2*x)});
\addplot[arkivblaa,thick,domain=0:2.5] ({-1.7*exp(-x)},{-1.7*exp(-2*x)});
\draw[->,black] (axis cs:1.2,.84706)--(axis cs:1,.58824);
\end{axis}
\begin{axis}[at={(nodeplot.east)},anchor=west,xshift=3mm,name=spiralplot,
 width=.31\textwidth,height=47mm,xmin=-2,xmax=2,ymin=-2,ymax=2,
 axis lines=middle,axis equal image,xlabel={$x_1$},ylabel={$x_2$},
 title={Stabil spiral},title style={font=\small},tick label style={font=\scriptsize},samples=180]
\addplot[arkivblaa,thick,domain=0:9] ({1.7*exp(-.35*x)*cos(deg(2*x))},{1.7*exp(-.35*x)*sin(deg(2*x))});
\addplot[arkivblaa,dashed,domain=0:9] ({-1.7*exp(-.35*x)*cos(deg(2*x))},{-1.7*exp(-.35*x)*sin(deg(2*x))});
\draw[->,black] (axis cs:1.7,0)--(axis cs:1.6088,.3261);
\end{axis}
\begin{axis}[at={(spiralplot.east)},anchor=west,xshift=3mm,
 width=.31\textwidth,height=47mm,xmin=-2,xmax=2,ymin=-2,ymax=2,
 axis lines=middle,axis equal image,xlabel={$x_1$},ylabel={$x_2$},
 title={Sadel},title style={font=\small},tick label style={font=\scriptsize},samples=80]
\addplot[arkivrod,thick,domain=-.6:.6] ({exp(x)},{exp(-x)});
\addplot[arkivrod,thick,domain=-.6:.6] ({-exp(x)},{exp(-x)});
\addplot[arkivrod,thick,domain=-.6:.6] ({exp(x)},{-exp(-x)});
\addplot[arkivrod,thick,domain=-.6:.6] ({-exp(x)},{-exp(-x)});
\draw[->,black] (axis cs:.819,1.221)--(axis cs:1.221,.819);
\end{axis}
\end{tikzpicture}
\caption{Fasebaner for tre lineære systemer $\boldsymbol{x}'=A\boldsymbol{x}$
i dimensionsløse koordinater. Fra venstre:
$A=\operatorname{diag}(-1,-2)$,
$A=\left(\begin{smallmatrix}-.35&-2\\2&-.35\end{smallmatrix}\right)$
og $A=\operatorname{diag}(1,-1)$. Pile angiver stigende tid.
Knude og spiral tiltrækker fremad i tiden; sadlen har både en stabil
og en ustabil retning. Akserne har samme skala i hvert panel.}
\label{fig:ode-fase}
\end{figure}


\subsection{Normvækst, energi og Lyapunovfunktioner}
Selv når alle egenværdier af en konstant matrix har negativ realdel,
kan en valgt norm vokse i en overgangsperiode, særligt for ikke-normale
matricer. Asymptotisk henfald betyder derfor ikke nødvendigvis
monotont henfald af hver komponent eller af enhver norm. En praktisk
eksponentialgrænse har formen
$\|\e^{At}\|\le M\e^{-\alpha t}$ for $t\ge0$, med $M\ge1$ og
$\alpha>0$ valgt efter matrixen og normen.

En Lyapunovfunktion er en skalar funktion, der kontrollerer afstanden
til en ligevægt uden at kræve en eksplicit løsning. For $V\in C^1$
gælder langs banen
$\dot V=\nabla V\cdot\boldsymbol{F}$. Lokal positiv definithed af $V$
og negativ definithed af $\dot V$ giver lokal asymptotisk stabilitet
under de almindelige lokale regularitetskrav. Et semidefinit fortegn
alene giver ikke samme konklusion; man skal undersøge den dynamik,
der kan blive i mængden $\dot V=0$.

For den udrivne normaliserede oscillator er
\[
 E=\tfrac12\bigl[(y')^2+\omega_0^2y^2\bigr],\qquad
 E'=-2\gamma(y')^2.
\]
Den fysiske energi er $mE$, hvis massen er $m$. For $\gamma=0$ er
$E$ bevaret; ved positiv dæmpning er den ikke-voksende. At $E'$ kan
være nul i isolerede vendepunkter er ikke en ny ligevægt. Henfaldet
til origo ved $\gamma>0$ følger også direkte af de homogene løsninger.

\subsection{Perturbation, tidsintervaller og numerisk kontrol}
En formel udvikling
$y(t;\varepsilon)=y_0(t)+\varepsilon y_1(t)+\cdots$ kræver, at parameteren
er gjort dimensionsløs, og at størrelsesordener sammenlignes på et
specificeret tidsinterval. Led, der vokser med tiden, kan ødelægge en
ellers lille parameterudvikling på lange tider. Et sekulært bidrag er
derfor en advarsel om manglende uniformitet, ikke nødvendigvis om en
ustabil fysisk løsning.

Hvis den lille parameter multiplicerer den højeste afledte, kan grænsen
sænke ligningens orden. Initial- eller randbetingelser kan da ikke alle
opfyldes af den reducerede ligning, og grænselag kan opstå. Dette er en
singulær perturbation; den kan ikke håndteres ved ukritisk at indsætte
en regulær potensrække. En mere omfattende perturbationsteori udvikles
senere efter behov.

Ved numerisk integration kontrolleres både lokal diskretiseringsfejl og
stabilitet. Stive systemer kan have hurtigt henfaldende moder, som tvinger
et eksplicit tidssteg til at være meget mindre end den langsomme fysiske
tidsskala. En lille residual er nyttig, men løsningsfejlen skal ses i
lyset af et følsomhedsestimat som \eqref{eq:ode-residual}. Kontrol af
bevarede størrelser eller dokumenteret energihenfald er et yderligere
fysisk konsistenstjek, ikke en erstatning for konvergenskontrol.

\section{Randbetingelser og metodeoverblik}
Initialdata i ét punkt og randbetingelser i flere punkter er forskellige
problemer. Entydighed for begyndelsesværdiproblemet giver ikke automatisk
entydighed for et randværdiproblem. Et homogent randproblem kan have
ikke-trivielle løsninger ved særlige spektralværdier; dette forbindes
senere med Sturm--Liouville-teori og Greenfunktioner.

Klassificer først orden, linearitet, autonomi og regulære intervaller.
Undersøg derefter, om ligningen er separabel, lineær, eksakt eller kan
transformeres til en kendt struktur. Ved lineære ligninger organiseres
løsningen som et homogent bidrag plus et partikulært bidrag. Ved
ikke-lineære systemer er kvalitativ analyse og lokale sætninger ofte
mere informative end et forsøg på en elementær formel.

En færdig løsningsbeskrivelse angiver interval, gren, alle mulige
ligevægte og de relevante initial-/randdata. Kontroller den i den
oprindelige ligning og kontroller dimensioner og fortegn. Ved brug af
egenværdier skal gentagne rødder, manglende egenvektorer og nulegenværdier
behandles særskilt. Konstante koefficienter, linearisering og små
parametre er forskellige antagelser og må ikke blandes stiltiende.

\chapter{Partielle differentialligninger}
\planlagt{Definitioner, notation, forudsætninger, centrale resultater og fysikforbindelser udvikles samlet.}
\section{Klassifikation og randbetingelser}
\leavevmode\par
% Redaktionelt: tilføj definitioner, gyldighed, relationer og relevante figurer.

\section{Separation af variable}
\leavevmode\par
% Redaktionelt: tilføj definitioner, gyldighed, relationer og relevante figurer.

\section{Varme-, bølge- og Laplaceligningen}
\leavevmode\par
% Redaktionelt: tilføj definitioner, gyldighed, relationer og relevante figurer.

\section{Karakteristikker og svage løsninger}
\leavevmode\par
% Redaktionelt: tilføj definitioner, gyldighed, relationer og relevante figurer.

\chapter{Konvolution, deltafunktion og Greenfunktioner}
\planlagt{Definitioner, notation, forudsætninger, centrale resultater og fysikforbindelser udvikles samlet.}
\section{Testfunktioner og distributioner}
\leavevmode\par
% Redaktionelt: tilføj definitioner, gyldighed, relationer og relevante figurer.

\section{Deltafunktionen og koordinatskift}
\leavevmode\par
% Redaktionelt: tilføj definitioner, gyldighed, relationer og relevante figurer.

\section{Konvolution og fundamentalløsninger}
\leavevmode\par
% Redaktionelt: tilføj definitioner, gyldighed, relationer og relevante figurer.

\section{Greenfunktioner, randbetingelser og kausalitet}
\leavevmode\par
% Redaktionelt: tilføj definitioner, gyldighed, relationer og relevante figurer.

\chapter{Sturm–Liouville-teori}
\planlagt{Definitioner, notation, forudsætninger, centrale resultater og fysikforbindelser udvikles samlet.}
\section{Selvadjungerede operatorer og vægtfunktion}
\leavevmode\par
% Redaktionelt: tilføj definitioner, gyldighed, relationer og relevante figurer.

\section{Randbetingelser og randled}
\leavevmode\par
% Redaktionelt: tilføj definitioner, gyldighed, relationer og relevante figurer.

\section{Ortogonalitet, spektrum og fuldstændighed}
\leavevmode\par
% Redaktionelt: tilføj definitioner, gyldighed, relationer og relevante figurer.

\section{Egenfunktionsudvikling og Greenkerner}
\leavevmode\par
% Redaktionelt: tilføj definitioner, gyldighed, relationer og relevante figurer.

\end{document}
